\documentclass[12pt,a4paper]{article}
\usepackage{amsmath,amssymb,amsthm}
\usepackage{geometry}
\geometry{top=3cm, bottom=3cm, left=2.5cm, right=2.5cm}
\usepackage{xcolor}
\usepackage{tcolorbox}
\tcbuselibrary{skins,breakable}

% تعریف رنگ‌های سفارشی برای زیبایی
\definecolor{primaryColor}{HTML}{2C3E50}
\definecolor{secondaryColor}{HTML}{18BC9C}
\definecolor{accentColor}{HTML}{E74C3C}

% تعریف کادرهای سوال و جواب و نکات کلیدی
\newtcolorbox{questionBox}[1]{
  colback=red!5!white,
  colframe=red!75!black,
  fonttitle=\bfseries,
  title={سوال احتمالی استاد: #1},
  fontupper=\small,
  breakable
}

\newtcolorbox{answerBox}[1]{
  colback=green!5!white,
  colframe=green!75!black,
  fonttitle=\bfseries,
  title={پاسخ پیشنهادی شما: #1},
  fontupper=\small,
  breakable
}

\newtcolorbox{tipBox}[1]{
  colback=blue!5!white,
  colframe=blue!75!black,
  fonttitle=\bfseries,
  title={نکته کلیدی برای ارائه: #1},
  fontupper=\small,
  breakable
}

\usepackage{xepersian}
\settextfont{Amiri}
\setlatintextfont{Times New Roman}

\title{\textbf{راهنمای جامع ارائه و دفاع شفاهی پروژه پوشش رأس وزن‌دار \lr{(Weighted Vertex Cover)}}}
\author{درس بهینه‌سازی ترکیبیاتی}
\date{\today}

\begin{document}

\maketitle

\begin{abstract}
\textbf{چکیده راهنما:} 
این سند به عنوان راهنمای کمکی دفاع شفاهی دانشجو برای ارائه پروژه پوشش رأس وزن‌دار \lr{(WVC)} تنظیم شده است. هدف این راهنما، هدایت گام‌به‌گام دانشجو در زمان ارائه پروژه به استاد، توضیح دلیل علمی کدهای پیاده‌سازی شده در نوت‌بوک، و نحوه پاسخ به سوالات سخت و چالش‌برانگیز داوران یا استاد درس است. مفاهیم پیاده‌سازی شده در قالب روابط ریاضی و تحلیل‌های رفتاری الگوریتم‌ها دسته‌بندی شده‌اند.
\end{abstract}

\newpage
\tableofcontents
\newpage

\section{ساختار فایل‌های خروجی پروژه}

پروژه نهایی تحویل داده شده شامل فایل‌های مهندسی‌شده زیر است که هر کدام وظیفه مشخصی دارند:
\begin{itemize}
    \item \textbf{\lr{Weighted\_Vertex\_Cover.ipynb}:} فایل اصلی نوت‌بوک شامل تعاریف، پیاده‌سازی کدهای آزمایش‌های عددی در ابعاد ۱۰۰ تا ۵۰۰ گره و نمودارهای رندر شده.
    \item \textbf{\lr{wvc\_tutorial.pdf} و \lr{wvc\_tutorial.tex}:} درسنامه علمی ۱۰ صفحه‌ای کامپایل شده به همراه اثبات‌های ریاضی و قضایا جهت مطالعه مبانی تئوری از پایه.
    \item \textbf{\lr{generate\_notebook.py} و \lr{execute\_notebook.py}:} اسکریپت‌های مکانیزه برای ایجاد و اجرای بدون خطای سلول‌های نوت‌بوک و ثبت خروجی‌ها به صورت متمرکز.
\end{itemize}

---

\section{سناریوی پیشنهادی برای ارائه گام‌به‌گام به استاد}

در زمان ارائه نوت‌بوک، بهتر است کدهای هر سلول را همراه با «چرایی» و «دلیل علمی» آن به استاد نشان دهید. از سناریوی زیر استفاده کنید:

\subsection{گام اول: تولید داده تصادفی گراف}
توضیح دهید که یک گراف تصادفی غیرجهت‌دار بر اساس مدل \lr{Erdős-Rényi} با چگالی یال $p=0.05$ تولید کرده و به هر گره یک وزن تصادفی بین ۱۰ تا ۱۰۰ اختصاص داده‌اید.
\begin{itemize}
    \item \textbf{چرا چگالی $0.05$؟} زیرا این مقدار تضمین می‌کند که گراف هم به اندازه کافی جاده (یال) برای پوشش دادن دارد و هم بیش از حد شلوغ نیست که مسئله بدیهی شود.
    \item \textbf{چرا وزن‌های مثبت؟} زیرا اگر وزن‌ها منفی یا صفر باشند، مسئله پوشش رأس وزن‌دار معنای فیزیکی و هزینه‌ای خود را از دست می‌دهد و گره‌های با وزن منفی همواره بدون هزینه انتخاب می‌شوند.
\end{itemize}

\subsection{گام دوم: حل با مدل برنامه‌ریزی اعداد صحیح باینری \lr{(IP)}}
با استفاده از کتابخانه \lr{PuLP} مدل را تعریف کرده و متغیرهای تصمیم‌گیری را به صورت باینری صفر و یک تعریف کرده‌اید.
\begin{tipBox}{محدودیت زمان حل \lr{IP}}
حتماً به استاد بگویید: «در زمان فراخوانی سالور \lr{CBC}، پارامتر محدودیت زمانی را روی ۱۵ ثانیه تنظیم کرده‌ام (\lr{timeLimit=15})؛ زیرا در ابعاد بالای ۱۵۰ رأس، مسئله به دلیل نمایی بودن فضای جستجوی درخت شاخه و حد ($2^N$) بسیار زمان‌بر می‌شود. این محدودیت زمانی باعث می‌شود سالور پس از ۱۵ ثانیه بهترین جواب شدنی پیدا شده را برگرداند تا روند اجرای کل پروژه متوقف نشود.»
\end{tipBox}

\subsection{گام سوم: آزادسازی خطی \lr{(LP Relaxation)} و بررسی حالت بدون سقف}
آزادسازی خطی را یک بار با قید استاندارد $0 \le x_i \le 1$ و بار دیگر با قید $x_i \ge 0$ (بدون حد بالای ۱) اجرا کرده‌اید.
\begin{itemize}
    \item \textbf{توضیح برای استاد:} به استاد نشان دهید که جواب بهینه هر دو مدل کاملاً یکسان است و بزرگترین مقدار متغیر در حالت بدون سقف نیز هرگز از ۱ تجاوز نکرده است.
    \item \textbf{علت تئوری:} چون وزن گره‌ها مثبت است ($w_i > 0$) و ما در حال کمینه کردن وزن کل هستیم، اگر در نقطه بهینه متغیر گرهی از ۱ بزرگتر باشد ($x_u^* > 1$)، سالور می‌تواند بدون نقض قیود یال‌ها ($1 + x_v^* \ge 1$) مقدار آن را به ۱ کاهش داده و مقدار تابع هدف را بهبود بخشد که این موضوع با فرض بهینه بودن در تناقض است.
\end{itemize}

\subsection{گام چهارم: الگوریتم گرد کردن \lr{(Rounding)}}
مقادیر اعشاری خروجی \lr{LP} را با آستانه $0.5$ گرد کرده و یک جواب باینری شدنی ساخته‌اید.
\begin{itemize}
    \item \textbf{اثبات شدنی بودن:} برای هر یال $(u, v)$ قید $x_u^* + x_v^* \ge 1$ برقرار است، پس حداقل یکی از این دو متغیر باید $\ge 0.5$ باشد (طبق اصل لانه کبوتر). بنابراین در زمان گرد کردن، حداقل یکی از آن‌ها به ۱ تبدیل شده و یال به صورت خودکار پوشش داده می‌شود.
    \item \textbf{ضریب تقریب ۲:} اثبات ریاضی نشان می‌دهد وزن جواب گرد شده همواره کمتر یا مساوی ۲ برابر کران پایین \lr{LP} و در نتیجه کمتر یا مساوی ۲ برابر بهینه واقعی \lr{IP} است.
\end{itemize}

\subsection{گام پنجم: بهبود به کمک جستجوی محلی \lr{(Local Search)}}
الگوریتم جستجوی محلی بر پایه روش \lr{Hill Climbing} را اجرا کرده‌اید که دو حرکت اصلی دارد: حذف گره‌های افزونه (که همسایگانشان در پوشش هستند) با ترتیب وزن نزولی، و حرکت تعویضی ۱-\lr{Swap} برای جایگزینی گره‌های سنگین با سبک‌تر.
\begin{itemize}
    \item \textbf{توضیح برای استاد:} جواب گرد شده در گراف‌های تصادفی شلوغ است (چون سالور تمام مقادیر را $0.5$ می‌گیرد و همه به ۱ گرد می‌شوند). جستجوی محلی با هزینه محاسباتی ناچیز (زیر ۱ میلی‌ثانیه)، این افزونگی‌ها را حذف کرده و جواب را به شدت به بهینه واقعی نزدیک می‌کند.
\end{itemize}

---

\section{آنالیز نتایج تجربی شبیه‌سازی}

استاد بر اساس خروجی‌های جدول آزمایش‌ها ممکن است چند تحلیل بخواهد:

\subsection{پدیده نسبت تقریب دقیقاً ۲.۰ در گرد کردن}
در جدول نتایج ابعاد بزرگ (۱۵۰ رأس به بالا)، متوجه می‌شوید که نسبت جواب گرد شده به \lr{LP} دقیقاً برابر با $2.0$ است.
\begin{itemize}
    \item \textbf{علت:} در گراف‌های تصادفی متقارن با چگالی متوازن و وزن‌های یکنواخت، بهینه‌ترین جواب برای آزادسازی خطی این است که به تمام متغیرها مقدار دقیقاً $0.5$ اختصاص داده شود ($x_i^* = 0.5$). در زمان گرد کردن، طبق قانون آستانه، تمام این مقادیر $0.5$ به ۱ تبدیل می‌شوند؛ یعنی کل گره‌های گراف در پوشش قرار می‌گیرند. این موضوع باعث می‌شود وزن پوشش گرد شده دقیقاً $\sum w_i$ شود که دقیقاً ۲ برابر جواب \lr{LP} (یعنی $0.5 \sum w_i$) است.
\end{itemize}

\subsection{کارایی فوق‌العاده جستجوی محلی}
با وجود اینکه جواب گرد شده بسیار سنگین است و نسبت تقریب ۲ را نشان می‌دهد، اما جستجوی محلی بلافاصله وارد عمل شده و گره‌های اضافی را حذف می‌کند. در جدول خروجی مشخص است که در سایز ۵۰۰ گره، وزن جواب جستجوی محلی با جواب بهینه \lr{IP} کمتر از $0.2\%$ تفاضل دارد، در حالی که زمان اجرای آن نزدیک به صفر است.

---

\section{سؤالات چالشی استاد در جلسه دفاع و نحوه پاسخ به آن‌ها}

\begin{questionBox}{چرا از ابتدا مستقیماً مدل دقیق \lr{IP} را برای ابعاد بزرگ حل نمی‌کنیم؟}
\textbf{پاسخ پیشنهادی:} زیرا مسئله پوشش رأس وزن‌دار یک مسئله \lr{NP-Hard} است. در ابعاد بزرگ، زمان حل درخت شاخه و حد به صورت نمایی با مرتبه زمانی $2^N$ رشد می‌کند و پیدا کردن جواب بهینه مطلق غیرممکن می‌شود. اما آزادسازی خطی \lr{LP} در زمان چندجمله‌ای (چند میلی‌ثانیه) حل می‌شود و به ما کمک می‌کند با روش گرد کردن و جستجوی محلی، در کسری از ثانیه به جوابی با اختلاف کمتر از ۳٪ نسبت به بهینه واقعی دست یابیم.
\end{questionBox}

\begin{questionBox}{در حذف گره‌های افزونه، چرا کاندیداها را بر اساس وزن نزولی مرتب کردید؟}
\textbf{پاسخ پیشنهادی:} هدف ما کمینه کردن وزن کل پوشش است. حذف یک گره افزونه ممکن است باعث شود همسایگان آن دیگر افزونه نباشند (یعنی نتوانیم آن‌ها را حذف کنیم). بنابراین، با مرتب‌سازی نزولی بر اساس وزن، مطمئن می‌شویم که ابتدا سنگین‌ترین گره‌های افزونه را حذف می‌کنیم تا بیشترین کاهش هزینه در تابع هدف حاصل شود.
\end{questionBox}

\begin{questionBox}{آیا جواب گرد شده همواره شدنی است؟ چطور این را اثبات می‌کنید؟}
\textbf{پاسخ پیشنهادی:} بله، ۱۰۰٪ شدنی است. برای هر یال $(u, v)$ قید $x_u^* + x_v^* \ge 1$ در مدل \lr{LP} ارضا می‌شود. طبق اصل لانه کبوتر، غیرممکن است هر دو متغیر کوچک‌تر از $0.5$ باشند؛ زیرا در این صورت جمعشان کوچک‌تر از ۱ خواهد شد. پس حداقل یکی از آن‌ها $\ge 0.5$ است و طبق آستانه به ۱ گرد می‌شود. در نتیجه، برای تمام یال‌ها حداقل یکی از دو سر یال در پوشش قرار می‌گیرد و جواب شدنی باقی می‌ماند.
\end{questionBox}

\begin{questionBox}{مفهوم شکاف بهینگی \lr{(Optimality Gap)} در آزمایش‌ها چیست؟}
\textbf{پاسخ پیشنهادی:} این شکاف نشان‌دهنده فاصله درصد بین کران پایین (وزن حاصل از \lr{LP Relaxation}) و جواب بهینه گسسته (\lr{IP}) است. هرچه این شکاف کوچک‌تر باشد، یعنی تخمین آزادسازی خطی به جواب واقعی نزدیک‌تر است. در آزمایش‌های تجربی ما این شکاف حول ۳۸٪ به دست آمد که نشان‌دهنده کارکرد طبیعی آزادسازی خطی در تخمین پوشش گره‌ها است.
\end{questionBox}

\end{document}
